%% 1-8 ながさを はかろう  第2部「はかる」の はじめの 章（第11回で 説明を ふやした）
%% 第1巻のきまり：分かち書き、漢字はすべてルビ、式は全角の＋－＝。数と単位のあいだは \,
%% この章の ものさし・テープの図は、本の上でも ほんとうの ながさ（目もり 1つ＝1 mm）で かいてある。
%% 子どもが じぶんの ものさしを あてて たしかめられる（印刷は 100％で）。

\makeatletter
%% ---- 1-8 の 絵の部品（\bkqg…） ----
% ほんとうの 大きさの ものさし  \bkqgruler{(x,y)}{さいごの cm}
%   (x,y) は 0 の 目もりの いちばん上。目もり 1つが 1 mm（本の上でも ほんものの 1 mm）
\newcommand\bkqgruler[2]{\begin{scope}[shift={#1}]
  \fill[bkYellowBg, draw=bkYellow!80!black, line width=0.8pt, rounded corners=1mm] (-4,-13) rectangle (#2*10+4,0);
  \pgfmathtruncatemacro\bkqg@n{#2*10}%
  \foreach \m in {0,...,\bkqg@n} {\pgfmathtruncatemacro\bkqg@c{mod(\m,10)==0}\pgfmathtruncatemacro\bkqg@h{mod(\m,5)==0}%
    \ifnum\bkqg@c=1 \draw[bkInk, line width=0.6pt] (\m,0) -- (\m,-5.4);
    \else\ifnum\bkqg@h=1 \draw[bkInk, line width=0.45pt] (\m,0) -- (\m,-3.8);
    \else\draw[bkInk, line width=0.3pt] (\m,0) -- (\m,-2.4);\fi\fi}
  \foreach \c in {0,...,#2} {\node[font=\sffamily\bfseries\footnotesize, text=bkInk, inner sep=0pt] at (\c*10,-9) {\c};}
\end{scope}}
% えんぴつ  \bkqgpencil{(x,y)}{ながさ}{いろ}   左はしが x、右はし（しんの\R{先}{さき}）が x＋ながさ。y が まん中（太さ 3 mm）
\newcommand\bkqgpencil[3]{\begin{scope}[shift={#1}]
  \fill[bkRed!55, rounded corners=0.4mm] (0,-1.5) rectangle (2,1.5);
  \fill[bkGray!45] (1.6,-1.5) rectangle (2.6,1.5);
  \fill[#3] (2.6,-1.5) rectangle (#2-4,1.5);
  \fill[#3!70!black] (2.6,-1.5) rectangle (#2-4,-0.5);
  \fill[bkSkin] (#2-4,-1.5) -- (#2-1.2,-0.4) -- (#2-1.2,0.4) -- (#2-4,1.5) -- cycle;
  \fill[bkInk!85] (#2-1.5,-0.52) -- (#2,0) -- (#2-1.5,0.52) -- cycle;
\end{scope}}
% けしゴム  \bkqgeraser{(x,y)}{ながさ}{カバーの いろ}   左下が (x,y)。上から 見た 形で、たかさ（はば）は ながさの 1/2.3
\newcommand\bkqgeraser[3]{\bkeraser[#3]{#1}{#2}{#2/2.3}}% 第11回：全巻共通の消しゴム（bk-pict1.sty）。本物に近い 2.3：1
% クレヨン  \bkqgcrayon{(x,y)}{ながさ}{いろ}   左はしが x、右はしが x＋ながさ。y が まん中（太さ 4 mm）
\newcommand\bkqgcrayon[3]{\begin{scope}[shift={#1}]
  \fill[#3, rounded corners=0.5mm] (0,-2) rectangle (#2-3.2,2);
  \fill[#3!22] (2.2,-2) rectangle (#2-6.5,2);
  \fill[#3!70] (2.2,-2) rectangle (#2-6.5,-1.3) (2.2,1.3) rectangle (#2-6.5,2);
  \fill[#3] (#2-3.4,-2) -- (#2,-0.8) -- (#2,0.8) -- (#2-3.4,2) -- cycle;
\end{scope}}
% 手のひら（上から 見た 手。ゆびは 右むき）  \bkqghand{(x,y)}{ながさ}{いろ}   手くびが x、ゆびの \R{先}{さき}が x＋ながさ
\newcommand\bkqg@handparts[1]{%
  \path[#1, rounded corners=0.9mm] (0,-4.2) rectangle (11,4.2);
  \foreach \bkqg@y/\bkqg@l in {3.2/17.4,1.07/20,-1.07/19.2,-3.2/16.6} {\path[#1, rounded corners=0.9mm] (9,\bkqg@y-0.95) rectangle (\bkqg@l,\bkqg@y+0.95);}
  \path[#1, rounded corners=0.8mm] (4.6,3.2) -- (9.6,7.6) -- (11.2,6.3) -- (7.4,2.4) -- cycle;}
\newcommand\bkqghand[3]{\begin{scope}[shift={#1}, scale=#2/20]
  \bkqg@handparts{draw=#3!65!black, line width=1.2pt}
  \bkqg@handparts{fill=#3}
  \draw[#3!65!black, line width=0.4pt] (9.4,2.1) -- (11.5,2.1) (9.4,0) -- (11.5,0) (9.4,-2.1) -- (11.5,-2.1);
\end{scope}}
% 工場  \bkqgfactory{(x,y)}{いろ}{かべの 文字}   左下が (x,y)、はば 20 mm、たかさ 17 mm（えんとつまで）
\newcommand\bkqgfactory[3]{\begin{scope}[shift={#1}]
  \fill[#2!80!black] (15,10) rectangle (17.6,17);
  \fill[#2] (0,0) -- (0,10) -- (6.6,13.5) -- (6.6,10) -- (13.3,13.5) -- (13.3,10) -- (20,13.5) -- (20,0) -- cycle;
  \node[font=\sffamily\bfseries\footnotesize, text=white, inner sep=0pt] at (10,4.6) {#3};
\end{scope}}
% ねじ（よこから 見た ところ）  \bkqgbolt{(x,y)}{ながさ}{いろ}   あたまの 左はしが x、y が まん中（あたまの たかさ 6.4 mm）
\newcommand\bkqgbolt[3]{\begin{scope}[shift={#1}]
  \filldraw[fill=#3!35, draw=#3!70!black, line width=0.5pt] (2.6,-1.6) -- (#2-0.5,-1.6) -- (#2,-1.1) -- (#2,1.1) -- (#2-0.5,1.6) -- (2.6,1.6) -- cycle;
  \pgfmathtruncatemacro\bkqg@t{#2-5}%
  \foreach \bkqg@k in {0,...,\bkqg@t} {\draw[#3!70!black, line width=0.4pt] (3.4+\bkqg@k,-1.6) -- (3.9+\bkqg@k,1.6);}
  \filldraw[fill=#3!55, draw=#3!70!black, line width=0.5pt, rounded corners=0.3mm] (0,-3.2) rectangle (2.6,3.2);
  \draw[#3!70!black, line width=0.4pt] (0,-1.1) -- (2.6,-1.1) (0,1.1) -- (2.6,1.1);
\end{scope}}
% ナット（まえから 見た ところ）  \bkqgnut{(x,y)}{はんけい}{いろ}   (x,y) が まん中
\newcommand\bkqgnut[3]{\begin{scope}[shift={#1}]
  \filldraw[fill=#3!55, draw=#3!70!black, line width=0.5pt] (0:#2) -- (60:#2) -- (120:#2) -- (180:#2) -- (240:#2) -- (300:#2) -- cycle;
  \filldraw[fill=white, draw=#3!70!black, line width=0.5pt] (0,0) circle (#2*0.45);
  \draw[#3!70!black, line width=0.3pt] (0,0) circle (#2*0.32);
\end{scope}}
\makeatother

\bkchapter{1}{8}{ながさを はかろう}{\bklevel{1}}{ながさを はかろう}
\bkmichishirube{1-3,1-6}{1-11,1-16,2-10}

\begin{goal}
\begin{itemize}[label={\color{bkTheme}●}, leftmargin=1.4\zw, itemsep=0pt, topsep=0pt]
\item 「1\,cm の いくつ\R{分}{ぶん}」で、 ながさを いえる。
\item ものさしで、 cm と mm の \R{目}{め}もりを よめる。
\item ながさの たし\R{算}{ざん}・ひき\R{算}{ざん}が できる。
\end{itemize}
\end{goal}

\begin{hatena}
\noindent{\sffamily\bfseries えんぴつ \R{何本}{なん|ぼん}\R{分}{ぶん}？}
\begin{bkfigure}
\begin{tikzpicture}[x=1mm, y=1mm]
  \fill[white, rounded corners=3mm] (0,0) rectangle (122,43);
  % つくえの よこ
  \draw[bkGray, line width=0.8pt, {Stealth[length=5pt]}-{Stealth[length=5pt]}] (22,38.5) -- (102,38.5);
  \node[fill=white, font=\sffamily\bfseries\footnotesize, text=bkGray, inner sep=1.5pt] at (62,38.5) {つくえの よこ};
  % つくえ（上から 見た ところ）
  \fill[bkStem!22, rounded corners=1.5mm] (22,3) rectangle (102,34);
  \draw[bkStem!75, line width=1.1pt, rounded corners=1.5mm] (22,3) rectangle (102,34);
  % A：ながい えんぴつ 4本
  \foreach \i in {0,...,3} {\bkqgpencil{(22+\i*20,24)}{20}{bkBlue}}
  \foreach \i in {1,...,4} {\node[font=\sffamily\bfseries\footnotesize, text=bkBlue!75!black] at (12+\i*20,29.5) {\i};}
  % B：みじかい えんぴつ 6本
  \foreach \i in {0,...,5} {\bkqgpencil{(22+\i*13.3333,10)}{13.3333}{bkGreen}}
  \foreach \i in {1,...,6} {\node[font=\sffamily\bfseries\footnotesize, text=bkGreen!60!black] at (15.3333+\i*13.3333,15.5) {\i};}
  % A・B の 子ども
  \node[circle, fill=bkYellow, text=white, font=\sffamily\bfseries, minimum size=6.5mm, inner sep=0pt] at (5,31) {A};
  \bkkid{(13,18)}{14.5}{bkBlue}{bkHair}
  \node[circle, fill=bkYellow, text=white, font=\sffamily\bfseries, minimum size=6.5mm, inner sep=0pt] at (5,15) {B};
  \bkkid{(13,3)}{14.5}{bkGreen}{bkHair}
  % けっか
  \node[anchor=west, font=\sffamily\bfseries, text=bkBlue!75!black] at (104,24) {4\R{本}{ほん}\R{分}{ぶん}};
  \node[anchor=west, font=\sffamily\bfseries, text=bkGreen!60!black] at (104,10) {6\R{本}{ぽん}\R{分}{ぶん}};
\end{tikzpicture}
\end{bkfigure}
\noindent おなじ つくえ なのに、 A は 4\R{本}{ほん}\R{分}{ぶん}、 B は 6\R{本}{ぽん}\R{分}{ぶん}。 \R{人}{ひと}に よって ちがったら こまるね。 どう すれば、 だれが はかっても おなじ \R{数}{かず}に なるかな？
\end{hatena}

\clearpage
%% ---------------- p.79 やってみよう ----------------
\begin{yattemiyou}[つくえの よこを はかろう]
\begin{enumerate}[label={\bkcircled[bkGreen]{\arabic*}}, leftmargin=2.2\zw, itemsep=2pt, topsep=0pt]
\item じぶんの えんぴつで はかろう。 はしを あわせて、 ずらしながら かぞえるよ。
\end{enumerate}
\begin{bkfigure}
\begin{tikzpicture}[x=1mm, y=1mm]
  \fill[bkStem!22] (8,0) rectangle (108,16);
  \draw[bkStem!75, line width=1.1pt] (8,16) -- (8,0) -- (108,0) -- (108,16);
  \foreach \i in {0,...,3} {\bkqgpencil{(8+\i*22,8)}{22}{bkOrange}}
  \foreach \i in {1,...,4} {\node[font=\sffamily\bfseries\footnotesize, text=bkOrange!80!black] at (-3+\i*22,13) {\i};}
  \draw[bkGreen!70!black, line width=0.9pt, {Stealth[length=3pt]}-{Stealth[length=3pt]}] (97,8) -- (107.6,8);
  \node[anchor=south, font=\sffamily\bfseries\footnotesize, text=bkGreen!70!black] at (102,9.5) {すこし};
  \node[anchor=north, font=\sffamily\bfseries\small, text=bkGreen!70!black] at (58,-1) {4\R{本}{ほん}と すこし};
\end{tikzpicture}
\end{bkfigure}
\begin{enumerate}[label={\bkcircled[bkGreen]{\arabic*}}, leftmargin=2.2\zw, itemsep=2pt, topsep=0pt, start=2]
\item こんどは、 \R{手}{て}のひらを ならべて はかろう。
\end{enumerate}
\begin{bkfigure}
\begin{tikzpicture}[x=1mm, y=1mm]
  \fill[bkStem!22] (8,0) rectangle (108,16);
  \draw[bkStem!75, line width=1.1pt] (8,16) -- (8,0) -- (108,0) -- (108,16);
  \foreach \i in {0,...,3} {\bkqghand{(8.4+\i*22.5,7.5)}{21.6}{bkSkin}}
  \foreach \i in {1,...,4} {\node[circle, fill=white, font=\sffamily\bfseries\footnotesize, text=bkOrange!80!black, inner sep=0.6pt] at (-5.6+\i*22.5,7.5) {\i};}
  \draw[bkGreen!70!black, line width=0.9pt, {Stealth[length=3pt]}-{Stealth[length=3pt]}] (98.6,8) -- (107.6,8);
  \node[anchor=south, font=\sffamily\bfseries\footnotesize, text=bkGreen!70!black] at (103.1,9.5) {すこし};
  \node[anchor=north, font=\sffamily\bfseries\small, text=bkGreen!70!black] at (58,-1) {4こと すこし};
\end{tikzpicture}
\end{bkfigure}
\begin{enumerate}[label={\bkcircled[bkGreen]{\arabic*}}, leftmargin=2.2\zw, itemsep=2pt, topsep=0pt, start=3]
\item \R{数}{かず}を かこう。 となりの \R{人}{ひと}と くらべよう。
\end{enumerate}
\par\vspace{4pt}
\noindent\hfill\begin{tblr}{colspec={Q[c,m,wd=30mm]Q[c,m,wd=40mm]Q[c,m,wd=40mm]}, hlines={0.6pt,bkLine}, vlines={0.6pt,bkLine},
  row{1}={bg=bkGreen!12, font=\sffamily\bfseries\small}, rowsep=5pt}
 & えんぴつ & \R{手}{て}のひら \\
じぶん & \bkkakikomi\enspace\R{本}{ほん}\R{分}{ぶん} & \bkkakikomi\enspace こ\R{分}{ぶん} \\
となりの \R{人}{ひと} & \bkkakikomi\enspace\R{本}{ほん}\R{分}{ぶん} & \bkkakikomi\enspace こ\R{分}{ぶん} \\
\end{tblr}\hfill\null
\end{yattemiyou}

\bkfukidashi{mina}{ソラと わたしで、 \R{数}{かず}が ちがったよ。 えんぴつも \R{手}{て}も、 \R{人}{ひと}に よって ながさが ちがうからだね。}

\clearpage
%% ---------------- 数学の道具（第11回：説明の段を ふやした） ----------------
\bksection{doguu}{\R{数学}{すう|がく}の \R{道具}{どう|ぐ}}

\bksubsection{1}{「1つ\R{分}{ぶん}」を そろえよう：cm}

\label{bk:1-8:cm}\bkhatenaref で、 A と B の \R{数}{かず}が ちがったのは なぜでしょう。 A の えんぴつは ながく、 B の えんぴつは みじかかったからです。 {\sffamily\bfseries\color{bkTheme}「1つ\R{分}{ぶん}」の ながさが ちがうと、 「いくつ\R{分}{ぶん}」も ちがって しまいます。}

\begin{bkfigure}
\begin{tikzpicture}[x=1mm, y=1mm]
  \fill[bkStem!22] (0,0) rectangle (60,16);
  \draw[bkStem!75, line width=1pt] (0,16) -- (0,0) -- (60,0) -- (60,16);
  \foreach \i in {0,...,3} {\bkqgpencil{(\i*15,11.5)}{15}{bkBlue}}
  \foreach \i in {0,...,5} {\bkqgpencil{(\i*10,4.5)}{10}{bkGreen}}
  \node[anchor=west, font=\sffamily\bfseries\footnotesize, text=bkBlue!75!black] at (63,11.5) {A：ながい えんぴつで 4\R{本}{ほん}\R{分}{ぶん}};
  \node[anchor=west, font=\sffamily\bfseries\footnotesize, text=bkGreen!60!black] at (63,4.5) {B：みじかい えんぴつで 6\R{本}{ぽん}\R{分}{ぶん}};
\end{tikzpicture}
\end{bkfigure}
\noindent そこで、 せかいの おおくの \R{国}{くに}で、 おなじ ながさの 「1つ\R{分}{ぶん}」、 1\,cm（センチメートル）を つかう ことに しました。\bkidx{せんちめーとる}{センチメートル（cm）} この 「1つ\R{分}{ぶん}」を、 たんい と いいます。\bkidx{たんい}{たんい}\bkidx{ながさ}{ながさ}

\begin{bkfigure}
\begin{tikzpicture}[x=1mm, y=1mm]
  % 1 cm
  \filldraw[fill=bkBlue!30, draw=bkBlue!80!black, line width=0.6pt] (0,0) rectangle (10,6);
  \node[font=\sffamily\bfseries\small, text=bkBlue!75!black] at (5,11) {1\,cm};
  % 5 cm ＝ 1 cm の 5つ分
  \foreach \i [count=\k] in {0,...,4} {\filldraw[fill=bkBlue!30, draw=bkBlue!80!black, line width=0.6pt] (20+\i*10,0) rectangle ++(10,6);
    \node[font=\sffamily\bfseries\footnotesize, text=bkBlue!75!black] at (25+\i*10,3) {\k};}
  \bkbrace{(20,7)}{(70,7)}{5\,cm}{1}
  \node[anchor=north, font=\sffamily\bfseries\footnotesize, text=bkBlue!75!black] at (45,-1) {1\,cm の 5つ\R{分}{ぶん}};
  \node[anchor=north, font=\sffamily\footnotesize, text=bkGray, align=center] at (5,-1) {ものさしを\\あてて みよう};
  % 5 cm の かきかた
  \begin{scope}[shift={(81,0)}]
    \fill[bkTheme!6, rounded corners=2.5mm] (0,-9) rectangle (45,17);
    \node[anchor=base east, font=\fontsize{30}{30}\selectfont, text=bkBlue!80!black, inner sep=0pt] at (15,3) {5};
    \node[anchor=base west, font=\fontsize{30}{30}\selectfont, text=bkOrange!90!black, inner sep=0pt] at (22.5,3) {cm};
    \node[font=\sffamily\bfseries\footnotesize, text=bkBlue!80!black] at (11,-4.5) {いくつ\R{分}{ぶん}};
    \node[font=\sffamily\bfseries\footnotesize, text=bkOrange!90!black] at (30.5,-4.5) {たんい};
  \end{scope}
\end{tikzpicture}
\end{bkfigure}
\noindent 1\,cm の 5つ\R{分}{ぶん}の ながさを、 5\,cm と かきます。 5は 「いくつ\R{分}{ぶん}」、 cm は 「なにの 1つ\R{分}{ぶん}か」を あらわします。\par
\noindent\bkhatenaref の こたえ：おなじ たんいで はかれば、 だれが はかっても おなじ \R{数}{かず}に なります。 \bkhatenaref の つくえは、 1\,cm の 60こ\R{分}{ぶん}で 60\,cm でした。 \R{学校}{がっ|こう}の つくえの よこは、 だいたい この ながさです。

\clearpage
\bksubsection{2}{ものさしで はかろう：cm と mm}

\label{bk:1-8:yomu}ものさしの \R{目}{め}もりを \R{大}{おお}きく して \R{見}{み}て みよう。 \R{目}{め}もりの \R{線}{せん}には、 ながさが 3しゅるい あります。

\begin{bkfigure}
\begin{tikzpicture}[x=1mm, y=1mm]
  % 0〜2 cm を 4ばいに して かく（1 mm → 4 mm）
  \fill[bkYellowBg, draw=bkYellow!80!black, line width=0.8pt, rounded corners=1.5mm] (-5,-24) rectangle (85,0);
  \foreach \m in {0,...,20} {\pgfmathtruncatemacro\c{mod(\m,10)==0}\pgfmathtruncatemacro\h{mod(\m,5)==0}%
    \ifnum\c=1 \draw[bkInk, line width=1.2pt] (\m*4,0) -- (\m*4,-12);
    \else\ifnum\h=1 \draw[bkInk, line width=0.9pt] (\m*4,0) -- (\m*4,-8);
    \else\draw[bkInk, line width=0.6pt] (\m*4,0) -- (\m*4,-4.5);\fi\fi}
  \foreach \c in {0,1,2} {\node[font=\sffamily\bfseries\small, text=bkInk] at (\c*40,-17.5) {\c};}
  \draw[bkTheme, line width=0.8pt, -{Stealth[length=3pt]}] (88,-6) -- (81.5,-10);
  \node[anchor=west, font=\sffamily\bfseries\footnotesize, text=bkTheme] at (88,-6) {ながい \R{線}{せん}：cm};
  \draw[bkBlue, line width=0.8pt, -{Stealth[length=3pt]}] (88,-14) -- (61.5,-7);
  \node[anchor=west, font=\sffamily\bfseries\footnotesize, text=bkBlue!80!black] at (88,-14) {まん\R{中}{なか}：5\,mm};
  \draw[bkGreen!70!black, line width=0.8pt, -{Stealth[length=3pt]}] (88,-21) -- (65.5,-4.5);
  \node[anchor=west, font=\sffamily\bfseries\footnotesize, text=bkGreen!60!black] at (88,-21) {みじかい \R{線}{せん}：1\,mm};
  \foreach \k in {1,...,10} {\node[font=\sffamily\bfseries\scriptsize, text=bkOrange!90!black] at (\k*4-2,3) {\k};}
\end{tikzpicture}
\end{bkfigure}
\begin{itemize}[label={\color{bkTheme}●}, leftmargin=1.4\zw, itemsep=0pt, topsep=0pt]
\item ながい \R{線}{せん}と ながい \R{線}{せん}の あいだが 1\,cm。 \R{線}{せん}の \R{下}{した}の 1、 2 は、 0から なん cm かを あらわす。
\item 1\,cm の あいだは、 10に \R{分}{わ}かれて いる。 いちばん みじかい あいだ 1つ\R{分}{ぶん}が 1\,mm（ミリメートル）。\bkidx{みりめーとる}{ミリメートル（mm）}
\item まん\R{中}{なか}の すこし ながい \R{線}{せん}は、 5\,mm の しるし。 かぞえやすく する ための \R{線}{せん}。
\end{itemize}

\noindent{\sffamily\bfseries\color{bkTheme}なぜ mm が いるの？}\quad えんぴつの ながさが、 ちょうど なん cm に ならない ときが あります。 1\,cm より みじかい 「はした」を あらわす ために、 1\,cm を 10に \R{分}{わ}けた mm を つかいます。

\clearpage
\noindent{\sffamily\bfseries\color{bkTheme}えんぴつの ながさを よもう}\par
\noindent ものさしの \R{目}{め}もりは、 ほんとうの ながさで かいて あります。 じぶんの ものさしを あてて みよう。

\begin{bkfigure}
\begin{tikzpicture}[x=1mm, y=1mm]
  \bkqgruler{(0,0)}{8}
  \bkqgpencil{(0,1.6)}{74}{bkBlue}
  \draw[bkTheme, line width=1pt, -{Stealth[length=4pt]}] (0,10.5) -- (0,4);
  \node[anchor=south west, font=\sffamily\bfseries\footnotesize, text=bkTheme, inner sep=1pt] at (-3,10.5) {0に あわせる};
  % 大きく した ところ（4ばい）
  \draw[bkGray, line width=0.6pt] (74,1) circle (4);
  \draw[bkGray, line width=0.6pt] (76.8,3.9) -- (88,17) (76.8,-1.9) -- (88,-11);
  \begin{scope}
    \clip[rounded corners=2.5mm] (88,-11) rectangle (128,17);
    \fill[white] (88,-11) rectangle (128,17);
    \fill[bkYellowBg] (88,-11) rectangle (128,2);
    \foreach \m in {70,...,78} {\pgfmathsetmacro\X{94+(\m-70)*4}%
      \ifnum\m=70 \draw[bkInk, line width=1.3pt] (\X,2) -- (\X,-20);
      \else\ifnum\m=75 \draw[bkInk, line width=1pt] (\X,2) -- (\X,-12.4);
      \else\draw[bkInk, line width=0.8pt] (\X,2) -- (\X,-6.8);\fi\fi}
    \fill[bkBlue] (88,2.4) rectangle (94,14.4);
    \fill[bkBlue!70!black] (88,2.4) rectangle (94,6.4);
    \fill[bkSkin] (94,2.4) -- (105.2,6.8) -- (105.2,10) -- (94,14.4) -- cycle;
    \fill[bkInk!85] (104,6.32) -- (110,8.4) -- (104,10.48) -- cycle;
    \foreach \k in {1,...,4} {\node[font=\sffamily\bfseries\footnotesize, text=bkOrange!90!black, inner sep=0pt] at (92+\k*4,-9) {\k};}
    \node[circle, fill=white, font=\sffamily\bfseries\small, text=bkInk, inner sep=0.5pt] at (90.8,-5) {7};
  \end{scope}
  \draw[bkGray, line width=0.6pt, rounded corners=2.5mm] (88,-11) rectangle (128,17);
\end{tikzpicture}
\end{bkfigure}
\begin{enumerate}[label={\bkcircled[bkTheme]{\arabic*}}, leftmargin=2.2\zw, itemsep=0pt, topsep=0pt]
\item ものさしの 0に、 えんぴつの はしを あわせる（1-3）。
\item えんぴつの \R{先}{さき}が、 どの \R{目}{め}もりに あるかを \R{見}{み}る。
\item \R{先}{さき}の \R{左}{ひだり}がわで、 いちばん \R{近}{ちか}い ながい \R{線}{せん}は 7。 0から 7まで、 1\,cm が 7つ。 だから 7\,cm。
\item 7の \R{線}{せん}から \R{先}{さき}まで、 みじかい あいだが 4つ。 だから 4\,mm。
\item あわせて、 えんぴつは 7\,cm 4\,mm。
\end{enumerate}
\noindent{\sffamily\bfseries\color{bkTheme}5\,mm の \R{線}{せん}を つかうと はやい}\quad \R{先}{さき}が 5の \R{線}{せん}より \R{右}{みぎ}なら、 「5 と あと いくつ」と かぞえます。 7\,cm 8\,mm なら、 5\,mm の \R{線}{せん}から 3つで 8\,mm です。

\clearpage
\bksubsection{3}{cm と mm の かんけい}

\label{bk:1-8:cmmm}1\,cm は 1\,mm が 10こ。 これは、 ばらが 10こで ぼうが 1\R{本}{ぽん}（1-2）と おなじ しくみです。

\begin{bkfigure}
\begin{tikzpicture}[x=1mm, y=1mm]
  \bkqgruler{(0,0)}{7}
  \foreach \c in {0,...,5} {\draw[bkBlue!80!black, line width=0.7pt, decorate, decoration={brace, amplitude=1.4mm}] (\c*10+0.4,1.6) -- (\c*10+9.6,1.6);
    \node[font=\sffamily\bfseries\scriptsize, text=bkBlue!80!black] at (\c*10+5,6.6) {10};}
  \draw[bkOrange!90!black, line width=0.7pt, decorate, decoration={brace, amplitude=1.2mm}] (60.4,1.6) -- (63.6,1.6);
  \node[font=\sffamily\bfseries\scriptsize, text=bkOrange!90!black] at (62,6.6) {4};
  \fill[bkTheme!70] (0,-13) rectangle (64,-15.5);
  \node[anchor=west, font=\sffamily\bfseries\footnotesize, text=bkInk] at (72,4) {10\,mm が 6こで 60\,mm};
  \node[anchor=west, font=\sffamily\bfseries\footnotesize, text=bkInk] at (72,-3) {60\,mm と 4\,mm で 64\,mm};
  \node[anchor=west, font=\sffamily\bfseries\footnotesize, text=bkTheme] at (72,-14.2) {6\,cm 4\,mm ＝ 64\,mm};
\end{tikzpicture}
\end{bkfigure}
\begin{enumerate}[label={\bkcircled[bkTheme]{\arabic*}}, leftmargin=2.2\zw, itemsep=0pt, topsep=0pt]
\item 6\,cm は、 1\,cm が 6こ。 1\,cm は 10\,mm だから、 10\,mm が 6こで 60\,mm（1-7の \R{何十}{なん|じゅう}）。
\item 60\,mm と 4\,mm を あわせて 64\,mm。
\end{enumerate}
\bkbigshiki{1\,cm ＝ 10\,mm}{おなじ ながさ}
\noindent{\sffamily\bfseries\color{bkTheme}ぎゃくに、 mm を cm と mm に}\par
75\,mm の テープを、 10\,mm ずつ いろを かえて \R{分}{わ}けて みよう。

\begin{bkfigure}
\begin{tikzpicture}[x=1mm, y=1mm]
  \foreach \i in {0,...,6} {\pgfmathtruncatemacro\c{mod(\i,2)}%
    \ifnum\c=0 \fill[bkBlue!45] (\i*10,0) rectangle ++(10,6);\else\fill[bkBlue!20] (\i*10,0) rectangle ++(10,6);\fi
    \node[font=\sffamily\bfseries\scriptsize, text=bkInk] at (\i*10+5,3) {10};}
  \fill[bkOrange!45] (70,0) rectangle (75,6);
  \node[font=\sffamily\bfseries\scriptsize, text=bkInk] at (72.5,3) {5};
  \draw[bkInk!60, line width=0.5pt] (0,0) rectangle (75,6);
  \bkbrace{(0,-1)}{(70,-1)}{10\,mm が 7こ ＝ 7\,cm}{-1}
  \node[anchor=west, font=\sffamily\bfseries\footnotesize, text=bkOrange!85!black] at (77,3) {のこり 5\,mm};
\end{tikzpicture}
\end{bkfigure}
\noindent 10\,mm の まとまりが 7こで 7\,cm、 のこりが 5\,mm。 だから 75\,mm ＝ 7\,cm 5\,mm です。 64\,mm なら、 まとまり 6こと のこり 4\,mm で 6\,cm 4\,mm。

\clearpage
\bksubsection{4}{ながい ものは m}

\label{bk:1-8:m}きょうしつの よこを cm で はかると、 900\,cm のように \R{大}{おお}きな \R{数}{かず}に なります。 そこで、 100\,cm を 1つ\R{分}{ぶん}に した 1\,m（メートル）を つかいます。\bkidx{めーとる}{メートル（m）}

\begin{bkfigure}
\begin{tikzpicture}[x=1mm, y=1mm]
  % 1 m ＝ 100 cm（10 cm ずつ ぬりわけ）
  \begin{scope}[shift={(0,8)}]
  \foreach \i in {0,...,9} {\pgfmathtruncatemacro\c{mod(\i,2)}%
    \ifnum\c=0 \fill[bkYellow!60] (\i*6,0) rectangle ++(6,5);\else\fill[bkYellow!25] (\i*6,0) rectangle ++(6,5);\fi}
  \draw[bkYellow!40!black, line width=0.6pt] (0,0) rectangle (60,5);
  \bkbrace{(0,6)}{(60,6)}{1\,m}{1}
  \foreach \x/\t in {0/0,30/50\,cm,60/100\,cm} {\node[anchor=north, font=\sffamily\bfseries\footnotesize, text=bkInk] at (\x,-0.5) {\t};}
  \node[anchor=north, font=\sffamily\footnotesize, text=bkGray] at (30,-6) {1つの いろが 10\,cm、 10こで 100\,cm};
  \end{scope}
  % きょうしつを 上から 見た ところ（1-21 と おなじ たて 7 m・よこ 9 m。1 m を 3.6 mm で かく）
  \begin{scope}[shift={(70,0)}]
    \fill[bkBlue!7] (0,0) rectangle (32.4,25.2);
    \draw[bkStem!75, line width=1pt] (0,0) rectangle (32.4,25.2);
    \fill[bkGreen!55!black, rounded corners=0.4mm] (7.2,23.1) rectangle (25.2,24.5);
    \fill[bkStem!50, rounded corners=0.3mm] (13.2,18.4) rectangle (19.2,20.8);
    \foreach \i in {0,...,3} {\foreach \j in {0,...,2} {\fill[bkStem!35, rounded corners=0.3mm] (4.2+\i*6.8,3+\j*4.8) rectangle ++(3.6,2.4);}}
    % よこ：1 m の ものさし 9本
    \foreach \i [count=\k] in {0,...,8} {\pgfmathtruncatemacro\c{mod(\i,2)}%
      \ifnum\c=0 \fill[bkYellow!60] (\i*3.6,-4.4) rectangle ++(3.6,3.6);\else\fill[bkYellow!25] (\i*3.6,-4.4) rectangle ++(3.6,3.6);\fi
      \node[font=\sffamily\bfseries\fontsize{7.5}{8}\selectfont, text=bkInk, inner sep=0pt] at (\i*3.6+1.8,-2.6) {\k};}
    \draw[bkYellow!40!black, line width=0.4pt] (0,-4.4) rectangle (32.4,-0.8);
    % たて：1 m の ものさし 7本
    \foreach \j [count=\k] in {0,...,6} {\pgfmathtruncatemacro\c{mod(\j,2)}%
      \ifnum\c=0 \fill[bkYellow!60] (33.2,\j*3.6) rectangle ++(3.6,3.6);\else\fill[bkYellow!25] (33.2,\j*3.6) rectangle ++(3.6,3.6);\fi
      \node[font=\sffamily\bfseries\fontsize{7.5}{8}\selectfont, text=bkInk, inner sep=0pt] at (35,\j*3.6+1.8) {\k};}
    \draw[bkYellow!40!black, line width=0.4pt] (33.2,0) rectangle (36.8,25.2);
    \node[anchor=north, font=\sffamily\bfseries\footnotesize, text=bkInk] at (16.2,-5.2) {きょうしつの よこ 9\,m};
    \node[anchor=west, font=\sffamily\bfseries\footnotesize, text=bkInk, align=left] at (37.8,12.6) {たて\\7\,m};
  \end{scope}
\end{tikzpicture}
\end{bkfigure}
\noindent きょうしつの よこは、 1\,m の ものさしで 9\R{本}{ほん}\R{分}{ぶん}。 だから 9\,m です。\\
1\,m と 20\,cm の ながさは、 1\,m 20\,cm と いいます。 1\,m は 100\,cm なので、 100\,cm と 20\,cm で 120\,cm と おなじです。
\bkbigshiki{1\,m ＝ 100\,cm}{おなじ ながさ}

\clearpage
\bksubsection{5}{ながさの 「かんかく」を もとう}

\label{bk:1-8:kankaku}みの まわりの ものの ながさを おぼえて おくと、 はかる \R{前}{まえ}に 「だいたい これくらい」と わかります。 はかった \R{数}{かず}が おかしい ときにも \R{気}{き}づけます。

\begin{bkfigure}
\begin{tikzpicture}[x=1mm, y=1mm]
  % 1円玉（ほんとうの 大きさ：はば 2 cm）
  \filldraw[fill=bkGray!18, draw=bkGray, line width=0.8pt] (10,22) circle (10);
  \draw[bkGray!70, line width=0.4pt] (10,22) circle (8.8);
  \node[font=\sffamily\bfseries\fontsize{18}{18}\selectfont, text=bkGray!90!black] at (10,22) {1};
  \draw[bkTheme, line width=0.8pt, {Stealth[length=3pt]}-{Stealth[length=3pt]}] (0,9.5) -- (20,9.5);
  \node[anchor=north, font=\sffamily\bfseries\footnotesize, text=bkTheme] at (10,8) {2\,cm};
  \node[anchor=north, font=\sffamily\footnotesize] at (10,3.5) {1\R{円玉}{えん|だま}};
  % はがき（よこ 10 cm、たて 14 cm 8 mm を 1/5 で かく）
  \begin{scope}[shift={(34,0)}]
  \filldraw[fill=white, draw=bkInk!60, line width=0.7pt] (0,8) rectangle (20,37.6);
  \fill[bkRed!60] (15,31) rectangle (18,35);
  \draw[bkTheme, line width=0.8pt, {Stealth[length=3pt]}-{Stealth[length=3pt]}] (0,5.5) -- (20,5.5);
  \node[anchor=north, font=\sffamily\bfseries\footnotesize, text=bkTheme] at (10,4) {10\,cm};
  \node[font=\sffamily\footnotesize] at (10,22) {はがき};
  \node[font=\sffamily\scriptsize, text=bkGray, align=center] at (10,14) {\R{小}{ちい}さく\\して ある};
  \end{scope}
  % この本（B5：よこ 18 cm 2 mm、たて 25 cm 7 mm を 1/7 で かく）
  \begin{scope}[shift={(74,0)}]
  \fill[bkTheme!15] (0,8) rectangle (26,44.7);
  \draw[bkTheme!70!black, line width=0.8pt] (0,8) rectangle (26,44.7);
  \node[font=\sffamily\bfseries\footnotesize, align=center] at (13,26) {この \R{本}{ほん}};
  \node[font=\sffamily\scriptsize, text=bkGray, align=center] at (13,17) {\R{小}{ちい}さく\\して ある};
  \draw[bkTheme, line width=0.8pt, {Stealth[length=3pt]}-{Stealth[length=3pt]}] (0,5.5) -- (26,5.5);
  \node[anchor=north, font=\sffamily\bfseries\footnotesize, text=bkTheme] at (13,4) {よこ 18\,cm 2\,mm};
  \draw[bkTheme, line width=0.8pt, {Stealth[length=3pt]}-{Stealth[length=3pt]}] (29,8) -- (29,44.7);
  \node[anchor=west, font=\sffamily\bfseries\footnotesize, text=bkTheme, align=left] at (30,26.4) {たて\\25\,cm\\7\,mm};
  \end{scope}
\end{tikzpicture}
\end{bkfigure}
\noindent 1\R{円玉}{えん|だま}は ほんとうの \R{大}{おお}きさで かいて あります。 はがきと この \R{本}{ほん}は \R{小}{ちい}さく して かいて あります。 ものさしで ほんものを はかって、 たしかめて みよう。

\noindent{\sffamily\bfseries\color{bkTheme}じぶんの からだの ものさし}\par
\noindent\hspace*{1\zw}\begin{tblr}{colspec={Q[l,m,wd=60mm]Q[l,m,wd=60mm]}, hlines={0.6pt,bkLine}, rowsep=4pt, colsep=4pt}
ゆびの はば（\R{人}{ひと}さしゆび） & だいたい\enspace\bkkakikomi\,cm \\
\R{手}{て}のひらの ながさ & だいたい\enspace\bkkakikomi\,cm \\
りょう\R{手}{て}を ひろげた ながさ & だいたい\enspace\bkkakikomi\,m\enspace\bkkakikomi\,cm \\
\end{tblr}
\par\vspace{2pt}
\noindent おうちの \R{人}{ひと}に 1\,m の ものさしや メジャーを もって もらって はかろう。 りょう\R{手}{て}を ひろげた ながさは、 じぶんの せの たかさと だいたい おなじです。

\clearpage
\bksubsection{6}{ながさの たし\R{算}{ざん}・ひき\R{算}{ざん}}

\label{bk:1-8:keisan}ながさも、 \R{数}{かず}と おなじように たしたり ひいたり できます。

\begin{bkfigure}
\begin{tikzpicture}[x=1mm, y=1mm]
  % つなぐ
  \node[anchor=west, font=\sffamily\bfseries\small, text=bkTheme] at (-1,15) {つなぐ → たし\R{算}{ざん}};
  \bktape{(0,0)}{30}{bkBlue}{3\,cm}
  \bktape{(30,0)}{40}{bkOrange}{4\,cm}
  \bkbrace{(0,-1)}{(70,-1)}{7\,cm}{-1}
  \node[font=\sffamily\bfseries\small, text=bkInk] at (35,-14) {3\,cm ＋ 4\,cm ＝ 7\,cm};
  % くらべる
  \begin{scope}[shift={(84,0)}]
  \node[anchor=west, font=\sffamily\bfseries\small, text=bkTheme] at (-1,15) {くらべる → ひき\R{算}{ざん}};
  \bktape{(0,2)}{40}{bkOrange}{4\,cm}
  \bktape{(0,-6)}{30}{bkBlue}{3\,cm}
  \draw[bkRed, line width=0.8pt, dashed] (30,-6) -- (30,9);
  \draw[bkRed, line width=1.2pt, {Stealth[length=3pt]}-{Stealth[length=3pt]}] (30.5,-2.5) -- (39.5,-2.5);
  \node[font=\sffamily\bfseries\footnotesize, text=bkRed!85!black] at (35,-8.5) {1\,cm};
  \node[font=\sffamily\bfseries\small, text=bkInk] at (20,-14) {4\,cm － 3\,cm ＝ 1\,cm};
  \end{scope}
\end{tikzpicture}
\end{bkfigure}
\noindent 2\R{本}{ほん}を つなぐと、 1\,cm が 3こと 4こで 7こ。 だから 7\,cm です。 くらべて ちがいを みる ときは、 ひき\R{算}{ざん}です（1-5）。

\noindent{\sffamily\bfseries\color{bkTheme}cm と mm が まじる とき}\par
cm は cm どうし、 mm は mm どうしで けいさんします。 くらいを そろえて たした ひっ\R{算}{さん}（1-7）と おなじです。
\begin{center}
4\,cm 2\,mm ＋ 3\,cm 5\,mm\quad →\quad cm：4 ＋ 3 ＝ 7、\enspace mm：2 ＋ 5 ＝ 7\quad →\quad 7\,cm 7\,mm
\end{center}
\begin{bkfigure}
\begin{tikzpicture}[x=1mm, y=1mm]
  \bkqgruler{(0,0)}{9}
  \fill[bkBlue!45] (0,1) rectangle (42,6);
  \fill[bkOrange!55] (42,1) rectangle (77,6);
  \node[font=\sffamily\bfseries\scriptsize, text=bkInk] at (21,3.5) {4\,cm 2\,mm};
  \node[font=\sffamily\bfseries\scriptsize, text=bkInk] at (59.5,3.5) {3\,cm 5\,mm};
  \draw[bkTheme, line width=1pt, -{Stealth[length=4pt]}] (77,12) -- (77,6.6);
  \node[anchor=south, font=\sffamily\bfseries\footnotesize, text=bkTheme] at (77,12) {7\,cm 7\,mm};
\end{tikzpicture}
\end{bkfigure}
\noindent ものさしの \R{上}{うえ}に ならべても、 ちょうど 7\,cm 7\,mm の ところで おわります。

\clearpage
\noindent{\sffamily\bfseries\color{bkTheme}mm が 10 より おおく なったら}\par
4\,cm 7\,mm ＋ 3\,cm 5\,mm は、 mm が 7 ＋ 5 ＝ 12で、 12\,mm に なります。

\begin{bkfigure}
\begin{tikzpicture}[x=1mm, y=1mm]
  \bkqgruler{(0,0)}{9}
  \fill[bkBlue!45] (0,1) rectangle (47,6);
  \fill[bkOrange!55] (47,1) rectangle (82,6);
  \node[font=\sffamily\bfseries\scriptsize, text=bkInk] at (23.5,3.5) {4\,cm 7\,mm};
  \node[font=\sffamily\bfseries\scriptsize, text=bkInk] at (64.5,3.5) {3\,cm 5\,mm};
  \draw[bkTheme, line width=1pt, -{Stealth[length=4pt]}] (82,12) -- (82,6.6);
  \node[anchor=south, font=\sffamily\bfseries\footnotesize, text=bkTheme] at (82,12) {8\,cm 2\,mm};
\end{tikzpicture}
\end{bkfigure}
\begin{enumerate}[label={\bkcircled[bkTheme]{\arabic*}}, leftmargin=2.2\zw, itemsep=0pt, topsep=0pt]
\item cm どうし：4 ＋ 3 ＝ 7。 mm どうし：7 ＋ 5 ＝ 12。
\item 12\,mm は、 10\,mm（＝1\,cm）と 2\,mm。
\item 1\,cm を cm の ほうへ くり\R{上}{あ}げる（1-7と おなじ）。 7 ＋ 1 ＝ 8で、 8\,cm 2\,mm。
\end{enumerate}
\noindent ものさしの \R{上}{うえ}でも、 ちょうど 8\,cm 2\,mm の ところで おわって います。

\begin{taisetsu}
ながさは、 たんいの いくつ\R{分}{ぶん}で あらわす。\\
1\,cm ＝ 10\,mm\qquad 1\,m ＝ 100\,cm\\
おなじ たんいどうしで けいさんする。
\end{taisetsu}

\bkkigocards{\bkkigocard{\sffamily cm}{センチメートル}{ながさの たんい。 1\,cm ＝ 10\,mm}{centimeter}{1-8}}
{\bkkigocard{\sffamily m}{メートル}{ながさの たんい。 1\,m ＝ 100\,cm}{meter}{1-8}}

\clearpage
%% ---------------- 例題 ----------------

\begin{reidai}[ものさしで よむ]
クレヨンの ながさは、 なん cm なん mm？
\begin{bkfigure}
\begin{tikzpicture}[x=1mm, y=1mm]
  \bkqgruler{(0,0)}{10}
  \bkqgcrayon{(0,2.1)}{63}{bkRed}
\end{tikzpicture}
\end{bkfigure}
\kangaekata はしは 0に あって いる。 \R{先}{さき}の \R{左}{ひだり}の ながい \R{線}{せん}は 6で、 6\,cm。 6の \R{線}{せん}から \R{先}{さき}まで、 みじかい あいだが 3つで 3\,mm。

\kotae 6\,cm 3\,mm
\end{reidai}

\begin{jibun}
けしゴムの ながさは、 なん cm なん mm？
\par\vspace{4pt}
\noindent\begin{minipage}[b]{80mm}
\begin{tikzpicture}[x=1mm, y=1mm]
  \bkqgruler{(4,0)}{7}
  \bkqgeraser{(4,0.3)}{45}{bkBlue}
\end{tikzpicture}
\end{minipage}\hfill
\begin{minipage}[b]{44mm}
{\sffamily\bfseries\small\color{bkOrange}\bkkotaelabel}\par\vspace{2pt}
\noindent\bkkakikomi\,cm\enspace\bkkakikomi\,mm
\end{minipage}
\par\vspace{3pt}
\noindent\begin{tikzpicture}[x=1mm, y=1mm]
  \bkqgruler{(4,0)}{9}
  \bkqgpencil{(4,1.6)}{86}{bkGreen}
\end{tikzpicture}\par\vspace{2pt}
\noindent\hfill えんぴつの ながさは？\enspace\bkkakikomi\,cm\enspace\bkkakikomi\,mm
\bkjans{けしゴム 4\,cm 5\,mm、 えんぴつ 8\,cm 6\,mm}
\end{jibun}

\begin{reidai}[cm と mm を かえる]
（1）6\,cm 4\,mm は、 なん mm？\hspace{2\zw}（2）75\,mm は、 なん cm なん mm？

\kangaekata （1）1\,cm ＝ 10\,mm だから、 6\,cm は 10\,mm が 6こで 60\,mm。 60\,mm と 4\,mm で 64\,mm。\\
（2）75\,mm は、 10\,mm の まとまりが 7こ（70\,mm）と、 のこり 5\,mm。 10\,mm が 7こで 7\,cm。

\kotae （1）64\,mm\hspace{3\zw}（2）7\,cm 5\,mm
\end{reidai}

\begin{jibun}
□に あう \R{数}{かず}を かきましょう。\par\vspace{3pt}
\noindent\hspace*{1\zw}（1）3\,cm 8\,mm ＝ \bkkakikomi[14mm]\,mm\hspace{3\zw}（2）47\,mm ＝ \bkkakikomi\,cm\enspace\bkkakikomi\,mm
\bkjans{（1）38\quad（2）4\,cm 7\,mm}
\end{jibun}

\begin{reidai}[m と cm を かえる]
（1）1\,m 30\,cm は、 なん cm？\hspace{2\zw}（2）150\,cm は、 なん m なん cm？

\kangaekata （1）1\,m ＝ 100\,cm。 100\,cm と 30\,cm で 130\,cm。\\
（2）150\,cm は、 100\,cm と 50\,cm。 100\,cm は 1\,m。

\kotae （1）130\,cm\hspace{3\zw}（2）1\,m 50\,cm
\end{reidai}

\begin{jibun}
□に あう \R{数}{かず}を かきましょう。\par\vspace{3pt}
\noindent\hspace*{1\zw}（1）1\,m 60\,cm ＝ \bkkakikomi[14mm]\,cm\hspace{3\zw}（2）120\,cm ＝ \bkkakikomi\,m\enspace\bkkakikomi\,cm
\bkjans{（1）160\quad（2）1\,m 20\,cm}
\end{jibun}

\begin{reidai}[つなぐ]
2\R{本}{ほん}の テープを つなぐと、 なん cm なん mm？
\par\vspace{2pt}\noindent\hfill\begin{tikzpicture}[x=1mm, y=1mm]
  \bktape{(0,0)}{42}{bkBlue}{4\,cm 2\,mm}
  \bktape{(42,0)}{35}{bkOrange}{3\,cm 5\,mm}
\end{tikzpicture}\hfill\null\par\vspace{2pt}
\kangaekata cm どうし 4 ＋ 3 ＝ 7、 mm どうし 2 ＋ 5 ＝ 7。 ものさしに ならべても、 7\,cm 7\,mm の ところで おわる。
\begin{bkfigure}
\begin{tikzpicture}[x=1mm, y=1mm]
  \bkqgruler{(0,0)}{9}
  \fill[bkBlue!45] (0,1) rectangle (42,5.5);
  \fill[bkOrange!55] (42,1) rectangle (77,5.5);
  \draw[bkTheme, line width=1pt, -{Stealth[length=4pt]}] (77,11) -- (77,6.2);
  \node[anchor=west, font=\sffamily\bfseries\footnotesize, text=bkTheme] at (78,10) {7\,cm 7\,mm};
\end{tikzpicture}
\end{bkfigure}

\shiki 4\,cm 2\,mm ＋ 3\,cm 5\,mm ＝ 7\,cm 7\,mm \kotaenext 7\,cm 7\,mm
\end{reidai}

\begin{jibun}
6\,cm 1\,mm と 2\,cm 6\,mm の テープを つなぐと？\par
\noindent\hspace*{1\zw}6\,cm 1\,mm ＋ 2\,cm 6\,mm ＝ \bkkakikomi\,cm\enspace\bkkakikomi\,mm
\bkjans{8\,cm 7\,mm}
\end{jibun}

\begin{reidai}[ちがい]
あかい リボンは 10\,cm 5\,mm、 あおい リボンは 6\,cm。 ながさの ちがいは？

\kangaekata cm どうし 10 － 6 ＝ 4。 あおい リボンは 0\,mm なので、 mm は 5の まま。

\shiki 10\,cm 5\,mm － 6\,cm ＝ 4\,cm 5\,mm \kotaenext 4\,cm 5\,mm
\end{reidai}



\begin{reidai}[mm が 10を こえる たし\R{算}{ざん}]
5\,cm 8\,mm と 2\,cm 4\,mm の テープを つなぐと？

\kangaekata cm どうし 5 ＋ 2 ＝ 7。 mm どうし 8 ＋ 4 ＝ 12。\\
12\,mm は 1\,cm と 2\,mm。 1\,cm を くり\R{上}{あ}げて、 7 ＋ 1 ＝ 8。

\shiki 5\,cm 8\,mm ＋ 2\,cm 4\,mm ＝ 8\,cm 2\,mm \kotaenext 8\,cm 2\,mm
\end{reidai}

\begin{jibun}
\noindent（1）14\,cm 6\,mm － 9\,cm ＝ \bkkakikomi\,cm\enspace\bkkakikomi\,mm\par\vspace{3pt}
\noindent（2）3\,cm 6\,mm ＋ 2\,cm 9\,mm ＝ \bkkakikomi\,cm\enspace\bkkakikomi\,mm
\bkjans{（1）5\,cm 6\,mm\quad（2）6\,cm 5\,mm（mm は 6 ＋ 9 ＝ 15で、 1\,cm と 5\,mm）}
\end{jibun}

\clearpage
%% ---------------- p.84 まちがえやすいところ ----------------
\bksection{tsumazuki}{まちがえやすい ところ}

\begin{tsumazuki}{0に あわせて いない}
\noindent\begin{minipage}[c]{0.6\linewidth}
\begin{tikzpicture}[x=1mm, y=1mm]
  \bkqgruler{(0,0)}{6}
  \bkqgeraser{(10,0.1)}{40}{bkGreen}
\end{tikzpicture}
\end{minipage}\hfill
\begin{minipage}[c]{0.38\linewidth}
けしゴムの ながさは？\par
\bkbatsu\enspace 5\,cm\hspace{2\zw}\bkmaru\enspace 4\,cm
\end{minipage}\par\vspace{4pt}
はしは 1の ところ。 5 － 1 ＝ 4 で 4\,cm。 0に あわせて はかろう。
\end{tsumazuki}

\begin{tsumazuki}{1\,m は 10\,cm？}
\bkbatsu\enspace 1\,m ＝ 10\,cm\hspace{3\zw}\bkmaru\enspace 1\,m ＝ 100\,cm\par
1\,cm ＝ 10\,mm と まぜないように しよう。
\end{tsumazuki}

\begin{tsumazuki}{cm と mm の \R{数}{かず}を たす}
\bkbatsu\enspace 4\,cm ＋ 3\,mm ＝ 7\,cm\par
\bkmaru\enspace 4\,cm ＋ 3\,mm ＝ 4\,cm 3\,mm\par
4と 3は、 たんいが ちがうので たせません。
\end{tsumazuki}

\begin{column}{1\,m の はじまり}
むかしは、 \R{手}{て}や \R{足}{あし}で ながさを はかって いました。 でも、 \R{人}{ひと}に よって ちがうので こまります。 そこで、 いまから 230\R{年}{ねん}ほど まえ、 フランスで 1\,m が きめられました。
\end{column}

\clearpage
%% ---------------- p.85 つかわれているところ ----------------
\begin{butsuri}
1\,mm や 1\,m の ながさは、 せかいの どこでも おなじです。 だから、 こんな ことが できます。
\begin{bkfigure}
\begin{tikzpicture}[x=1mm, y=1mm]
  % ねじと ナットを つくる
  \fill[white, rounded corners=3mm] (0,0) rectangle (59,44);
  \bkqgfactory{(3,25)}{bkBlue}{A の \R{国}{くに}}
  \bkqgfactory{(36,25)}{bkOrange}{B の \R{国}{くに}}
  % A の ねじ（よこから）
  \bkqgbolt{(4.5,19.5)}{13}{bkBlue}
  \node[anchor=north, font=\sffamily\bfseries\fontsize{8}{9}\selectfont, text=bkBlue!75!black, inner sep=0pt] at (11,15.3) {ねじ M6};
  % B の ナット（まえから）
  \bkqgnut{(48.5,19.5)}{3.6}{bkOrange}
  \node[anchor=north, font=\sffamily\bfseries\fontsize{8}{9}\selectfont, text=bkOrange!80!black, inner sep=0pt] at (48.5,15.3) {ナット M6};
  \draw[bk arrow, bkGray, line width=1pt] (19,18) -- (23.5,12.5);
  \draw[bk arrow, bkGray, line width=1pt] (42,18.5) -- (36.5,12.5);
  % くみあわせた ところ
  \bkqgbolt{(14,6)}{22}{bkBlue}
  \filldraw[fill=bkOrange!55, draw=bkOrange!70!black, line width=0.5pt, rounded corners=0.3mm] (27,2.4) rectangle (31.4,9.6);
  \draw[bkOrange!70!black, line width=0.5pt] (27,4.4) -- (31.4,4.4) (27,7.6) -- (31.4,7.6);
  \node[anchor=west, font=\sffamily\bfseries\footnotesize, text=bkInk] at (38,6) {ぴったり};
  % 家の ずめん
  \begin{scope}[shift={(63,0)}]
  \fill[white, rounded corners=3mm] (0,0) rectangle (59,44);
  \fill[bkBlue!12] (5,4) rectangle (54,40);
  \draw[bkBlue!70!black, line width=1.4pt] (15,10) rectangle (47,32);
  \draw[bkBlue!70!black, line width=0.8pt] (31,10) -- (31,32) (15,21) -- (31,21);
  \draw[bkBlue!70!black, line width=0.5pt, {Stealth[length=3pt]}-{Stealth[length=3pt]}] (15,36) -- (47,36);
  \node[fill=bkBlue!12, font=\sffamily\bfseries\footnotesize, text=bkBlue!70!black, inner sep=0.8pt] at (31,36) {7280};
  \draw[bkBlue!70!black, line width=0.5pt, {Stealth[length=3pt]}-{Stealth[length=3pt]}] (10.5,10) -- (10.5,32);
  \node[fill=bkBlue!12, font=\sffamily\bfseries\footnotesize, text=bkBlue!70!black, inner sep=0.8pt, rotate=90] at (10.5,21) {5460};
  \node[font=\sffamily\bfseries\footnotesize, text=bkInk] at (35,6.5) {\R{家}{いえ}の ずめん};
  \end{scope}
\end{tikzpicture}
\end{bkfigure}
\begin{itemize}[label={\color{bkPurple}●}, leftmargin=1.4\zw, itemsep=0pt, topsep=0pt]
\item \R{車}{くるま}や きかいは、 たくさんの ねじで とめて あります。 ねじの \R{太}{ふと}さは mm で きめます。 「M6」の ねじは、 \R{太}{ふと}さが 6\,mm です。 ちがう \R{国}{くに}の \R{工場}{こう|じょう}で つくった ねじと ナットでも、 ぴったり はまります。
\item \R{家}{いえ}の ずめんの ながさは mm で かきます。 7280 は 7280\,mm。 だれが よんでも、 おなじ \R{大}{おお}きさの \R{家}{いえ}が できます。
\end{itemize}
\end{butsuri}

\bkfukidashi{delta}{きかいを つくる \R{工場}{こう|じょう}では、 1\,mm を 1000こに \R{分}{わ}けた ながさまで はかる ことも あるよ！}

\begin{bkfigure}
\begin{tikzpicture}[x=1mm, y=1mm]
  \draw[bkRoute7, line width=3pt] (10,0) -- (128,0);
  \node[anchor=west, inner sep=0pt] at (0,12) {\bkrouteicon[1.4em]{7}\hspace{0.3em}{\sffamily\bfseries\small\color{bkRoute7!80!black}「はかる」の \R{路線}{ろ|せん}}};
  \foreach \x/\c/\t/\s in {10/bkVol1/{1-8}/{ながさ},33.6/bkVol1/{1-9}/{かさ・おもさ},57.2/bkVol1/{1-11}/{\R{目}{め}もり},80.8/bkVol1/{1-16}/{\R{小数}{しょう|すう}},104.4/bkVol2/{2-6}/{ずれ},128/bkVol3/{3-14}/{ばらつき}} {
    \fill[white, draw=\c, line width=1.6pt] (\x,0) circle (2.6);
    \node[anchor=base, font=\sffamily\bfseries\footnotesize, text=\c!85!black] at (\x,-8.4) {\t};
    \node[anchor=base, font=\sffamily\footnotesize, text=bkInk] at (\x,-15.4) {\s};}
  \node[fill=bkTheme, text=white, rounded corners=1mm, font=\sffamily\bfseries\fontsize{8}{9}\selectfont, inner sep=1.5pt] at (10,5.5) {いま ここ};
\end{tikzpicture}
\end{bkfigure}

\clearpage
%% ---------------- p.86〜87 もんだい ----------------
\bkendproblems

\bkproblemgroup{きほん}{1}

\bkproblem 0から ア・イ・ウ までの ながさは？
\begin{bkfigure}
\begin{tikzpicture}[x=1mm, y=1mm]
  \bkqgruler{(0,0)}{10}
  \foreach \x/\t in {30/ア,65/イ,92/ウ} {\draw[bkTheme, line width=1.2pt, -{Stealth[length=4pt]}] (\x,6.5) -- (\x,1);
    \node[font=\sffamily\bfseries\small, text=bkTheme] at (\x,9.6) {\t};}
\end{tikzpicture}
\end{bkfigure}
\noindent\hspace*{1\zw}ア\enspace\bkkakikomi[28mm]\hspace{2\zw}イ\enspace\bkkakikomi[28mm]\hspace{2\zw}ウ\enspace\bkkakikomi[28mm]
\bkpans{\mbox{ア 3\,cm}\quad \mbox{イ 6\,cm 5\,mm}\quad \mbox{ウ 9\,cm 2\,mm}}

\bkproblem □に あう \R{数}{かず}を かきましょう。
\par
\noindent\begin{tblr}{colspec={Q[l,m,wd=71mm]Q[l,m,wd=68mm]}, colsep=0pt, rowsep=1pt}
（1）1\,cm ＝ \bkkakikomi[13mm]\,mm & （2）1\,m ＝ \bkkakikomi[13mm]\,cm \\
（3）5\,cm 3\,mm ＝ \bkkakikomi[13mm]\,mm & （4）1\,m 40\,cm ＝ \bkkakikomi[13mm]\,cm \\
（5）82\,mm ＝ \bkkakikomi[10mm]\,cm\,\bkkakikomi[10mm]\,mm & （6）180\,cm ＝ \bkkakikomi[10mm]\,m\,\bkkakikomi[10mm]\,cm \\
\end{tblr}
\bkpans{（1）10\quad（2）100\quad（3）53\quad（4）140\quad（5）8\,cm 2\,mm\quad（6）1\,m 80\,cm}

\bkproblem けいさんを しましょう。
\par
\noindent\hspace*{1\zw}\begin{tblr}{colspec={Q[l,m,wd=130mm]}, colsep=0pt, rowsep=1pt}
（1）8\,cm ＋ 6\,cm ＝ \bkkakikomi[13mm]\,cm \\
（2）13\,cm － 5\,cm ＝ \bkkakikomi[13mm]\,cm \\
（3）4\,cm 5\,mm ＋ 3\,cm 2\,mm ＝ \bkkakikomi\,cm\enspace\bkkakikomi\,mm \\
（4）9\,cm 7\,mm － 5\,cm ＝ \bkkakikomi\,cm\enspace\bkkakikomi\,mm \\
（5）2\,cm 7\,mm ＋ 4\,cm 5\,mm ＝ \bkkakikomi\,cm\enspace\bkkakikomi\,mm \\
\end{tblr}
\bkpans{（1）14\,cm\quad（2）8\,cm\quad（3）7\,cm 7\,mm\quad（4）4\,cm 7\,mm\quad（5）7\,cm 2\,mm}

\bkproblem 9\,cm と 6\,cm の テープを つなぐと、 なん cm？
\bkshikiran
\bkpans{9\,cm ＋ 6\,cm ＝ 15\,cm\quad 15\,cm}

\bkproblemgroup{チャレンジ}{1}

\bkproblem ソラの せの たかさは 1\,m 22\,cm、 ミナは 119\,cm です。 どちらが、 なん cm たかいですか。\par
{\small\color{bkOrange!58!black}ヒント：1\,m 22\,cm を、 cm だけで かいて みよう。 ちがいは、 119から 122まで かぞえたすと わかるよ（1-7）。}
\bkshikiran[11\zw]
\bkpans{1\,m 22\,cm ＝ 122\,cm、 122\,cm － 119\,cm ＝ 3\,cm\quad ソラが 3\,cm たかい（119から 122まで かぞえたすと 3）}

\bkproblem 1\,m の ものさしと 30\,cm の ものさしを、 まっすぐ つなぎました。 ぜんぶで なん cm ですか。
\bkshikiran
\bkpans{100\,cm ＋ 30\,cm ＝ 130\,cm\quad 130\,cm（1\,m 30\,cm）}

\begin{furikaeri}
\begin{checklist}
\item 「1\,cm の いくつ\R{分}{ぶん}」で ながさを いえる。\hfill{\footnotesize\sffamily\color{bkTheme}→\pageref{bk:1-8:cm}ページ}
\item ものさしで cm と mm を よめる。\hfill{\footnotesize\sffamily\color{bkTheme}→\pageref{bk:1-8:yomu}ページ}
\item ながさの たし\R{算}{ざん}・ひき\R{算}{ざん}が できる。\hfill{\footnotesize\sffamily\color{bkTheme}→\pageref{bk:1-8:keisan}ページ}
\end{checklist}
チェックが つかなかったら、 →の ページに もどろう。
\end{furikaeri}
\clearpage
